% !TEX root = ../main.tex

\begin{digest}

\noindent\textbf{Background and objectives.}
Modern computer, communication, and consumer-electronics manufacturing must handle small components, frequent product changes, and contact-rich operations. Conventional fixed-program automation performs well in stable high-volume production, but each new component or assembly sequence can require substantial reprogramming, fixture redesign, and system integration. This project developed a learning-based dual-arm system with a language-conditioned policy interface and validated one shared checkpoint under three fixed scene--instruction task configurations, covering hardware selection, teleoperated demonstration collection, packet-associated multi-view observations, model training, real-robot deployment, and physical evaluation.

The final task scope consisted of three bimanual operations. In the ring task, the left arm repositioned a can while the right arm grasped a handled black ring, placed it over the can, and released it. In the socket task, the right arm pushed a box containing a two-hole socket to the middle of the workspace while the left arm grasped a two-pin plug and inserted it. In the memory-module task, the left arm pulled and stabilized a motherboard while the right arm removed the memory module and placed it in a yellow tray. These tasks were selected to cover repositioning, grasping, coordinated bilateral roles, constrained insertion, extraction, transport, and placement.

\noindent\textbf{Requirements and concept selection.}
Customer requirements were translated into six engineering specifications through quality function deployment. The design-stage targets were planning success above 90\%, task-execution success above 70\%, component-detection accuracy above 90\%, pose-estimation error below 20~mm, failure-detection rate above 90\%, and recovery success above 70\%. Planning, component detection, failure detection, and automatic recovery each used an independent 25-experiment protocol, with separate test sets for ES5 and ES6. Task execution used 25 physical trials for each of the three tasks, and all completed tasks were checked against the ES4 positional tolerance.

The concept-generation process separated the system into robot platform, teleoperation, data collection, end-effector, and policy modules. Alternatives were screened for safety, interface compatibility, availability, manufacturability, and schedule feasibility. The selected system uses two Franka Research~3 robots, two GELLO leader devices, one overhead camera, two wrist cameras, custom compliant grippers, and the LeRobot implementation of $\pi_{0.5}$. One unified checkpoint is trained jointly across the three tasks to accept language instructions and generate bilateral actions.

\noindent\textbf{Final system design.}
The workcell contains two upright FR3 arms operating in a shared tabletop region. Each arm contributes seven joint positions and one gripper value, producing a 16-dimensional bilateral proprioceptive state. The action interface uses the corresponding 16-dimensional ordering: seven absolute left-arm joint targets, left-gripper width, seven absolute right-arm joint targets, and right-gripper width. One fixed Intel RealSense D435 camera provides an overhead view of the shared scene, while two Intel RealSense D405 cameras provide close-range left- and right-wrist views. The three RGB streams provide complementary global and local observations near grasp, contact, insertion, extraction, and placement regions.

Demonstrations were collected bilaterally with two GELLO leaders. The operator controlled the left and right follower arms while the recording system stored packet-associated multi-view observations, measured robot and gripper state, absolute action targets, task identity, and episode metadata. The same task instruction used during collection was retained for training and deployment. Before collection, the system checked leader zero positions, channel directions, gripper mappings, camera identities, packet freshness, and command limits. Episodes were accepted when all three image streams and the complete state/action record were present and the intended task was completed correctly.

\noindent\textbf{Dataset and model training.}
The final dataset contains 240 bilateral episodes: 80 for the ring task, 80 for the socket task, and 80 for the memory-module task. In total, the dataset contains 119,650 frames. Eight episodes from each task were assigned to validation using the same fixed episode indices, leaving 216 episodes for training and 24 for validation. The split was performed at the episode level to prevent neighbouring frames from the same physical trajectory from appearing in both sets. During training, the three task datasets were selected with equal probability so that the ring dataset's larger frame count did not dominate the optimization.

The final controller was initialized from the official \texttt{pi05\_base} model and trained through LeRobot~0.5.2 with project-specific modifications. The PaliGemma vision-language backbone remained frozen; only the action expert and the action/time projection modules were trained, and no LoRA or other parameter-efficient adapter was used. The model received three resized and normalized RGB images, a language instruction, and the current bilateral state. The valid 16-dimensional state and action vectors were padded to the model-side representation while preserving their physical ordering. Joint training ran for 100,000 optimizer steps on an NVIDIA RTX~A6000 with batch size~4 and AdamW optimization. Validation was performed every 5,000 steps. For checkpoint selection, the normalized flow-matching losses of the three tasks were averaged without frame-count weighting, giving every task equal influence. The minimum task-balanced macro loss, 0.036552, occurred at 65,000 steps, and this checkpoint was selected for physical deployment.

\noindent\textbf{Deployment and safety boundary.}
Physical deployment used the selected unified checkpoint for all three instructions. The policy received the selected instruction, packet-associated multi-view observations, and the 16-dimensional bilateral state. It predicted a 50-step continuous bimanual action chunk, and the executor consumed the first ten actions before acquiring an updated observation and requesting another prediction. Demonstration collection and training trajectories used 5~Hz, whereas deployment action execution used 15~Hz.

Before enabling motion, a dry run verified the checkpoint identity, camera ordering and freshness, state/action dimensions, normalization statistics, instruction, network communication, and finite model output. Robot-side checks enforced joint, gripper, workspace, timeout, protective-stop, emergency-stop, and manual-stop conditions. Recoverable task-level failures invoked the automatic recovery procedure.

\noindent\textbf{Physical validation results.}
The selected checkpoint was evaluated on an NVIDIA GeForce RTX~5070~Ti in 75 physical trials, with 25 trials per task. Before each trial, the objects were reset around the nominal arrangement, with initial displacement kept within approximately 5~cm for the ring and memory-module scenes and approximately 3~cm for the socket scene. A trial was counted as successful when the complete task-specific terminal condition was achieved.

The memory-module task succeeded in all 25 trials, producing an observed success rate of 100\%. The ring task succeeded in 21 of 25 trials, producing 84\%. The plug-insertion task succeeded in 13 of 25 trials, producing 52\%. Across all tasks, the system completed 59 of 75 trials successfully, giving an overall observed success rate of 78.7\%. The overall value exceeded the 70\% task-execution target, and the memory-module and ring tasks exceeded the target individually. However, the socket task did not meet the target, demonstrating that the aggregate result did not represent uniform performance across all three operations.

The independent specification-level tests produced planning success of 23/25 (92\%), component detection of 24/25 (96\%), failure detection of 23/25 (92\%), and automatic recovery of 23/25 (92\%). ES5 and ES6 used two separate groups of 25 experiments. All 59 completed tasks satisfied the ES4 positional tolerance. These results met the project-level ES1, ES3, ES4, ES5, and ES6 thresholds.

\noindent\textbf{Failure analysis and discussion.}
The memory-module result indicates that the selected arm roles, camera observations, compliant grippers, demonstration data, and unified checkpoint worked reliably for the evaluated motherboard configuration.

Three of the four ring failures occurred before secure transport. In two trials, the right gripper approached away from the handle and missed the grasp; in one trial, the ring base obstructed the handle approach. The remaining failure occurred during placement over the can. The dominant limitation was therefore handle acquisition rather than the initial can repositioning or the overall bilateral sequence. Additional demonstrations around displaced handle poses, improved wrist-camera visibility, a more repeatable pre-grasp approach, or a revised finger profile could directly address this failure pattern.

The socket task exposed the most important remaining performance bottleneck. Two of its twelve failures occurred because the left gripper did not acquire or retain the plug. The other ten failures occurred after successful acquisition because the two pins were not aligned accurately enough with the two holes. Likely contributors include residual socket-box pose variation after pushing, plug-pose variation inside the compliant gripper, limited close-range visual information, and the interval over which ten predicted actions were executed before replanning.

\noindent\textbf{Conclusions and recommendations.}
The project produced a functioning and traceable VLA-based proof-of-concept system for unified multi-task bimanual manipulation. Its principal contribution is the integration of bilateral teleoperation, packet-associated multi-view observations, reproducible dataset construction, joint multi-task $\pi_{0.5}$ training, a unified 16-dimensional action interface, robot-side safeguards, and real-robot validation. The same checkpoint successfully executed all three fixed scene--instruction configurations. Because scene and instruction were not varied independently, the trials do not isolate the contribution of language or establish sensitivity to instruction variation.

Future development should first improve plug alignment and insertion through a more constrained fixture, better close-range visibility, and targeted demonstrations. The next system-level step is safe-space estimation from packet-associated multi-view observations, followed by broader object, scene, and instruction variation.

\end{digest}
